\section{The prescribed-factor extension and the point requiring repair}
\label{sec:diagnosis}

This section records the precise relation between the published
smooth-modulus theorem \cite{Stadlmann2025}, the S.~Type-I extension proposed
in \cite{Stadlmann2026}, and \cref{thm:S52}.  The distinction matters
because the larger modulus family is not covered by the published theorem.

\subsection{What the published theorem supplies}

The 2025 theorem assumes that the relevant modulus divides $P(x^\delta)$;
in particular it is squarefree and every prime factor is at most $x^\delta$.
Its q-van der Corput argument introduces a divisor at the scale
\[
 D=\frac{N}{x^{50\eta}H^2}.
\]
This theorem and its local reduction are established analytic inputs in the
present paper.  What must be proved afresh is that the reduction survives when
global smoothness is replaced by the two divisor witnesses in the definition
of $\DS$.

\subsection{The window inconsistency}

In Lemma~6 on page~12 of version~1 of \cite{Stadlmann2026}, the
S.~Type-I statement uses
\[
 r^2x^{2-\gamma-52\eta-\delta}d^{-4}<d_1<
 r^2x^{2-\gamma-52\eta}d^{-4}.
\]
At the beginning of its proof on page~17 this is restated with $100\eta$,
while the factor family immediately following the restatement again uses
$52\eta$.
These three displayed conditions define different modulus families.  The
proof therefore does not, as written, establish the stated lemma.

The compatible exponent is $52$.  Indeed, after the first factorization,
\[
 H=\frac{x^\eta RQ^2}{q_0M},
\]
and the stated divisor window gives
\[
 \frac{x^{2-\gamma-52\eta}}{Q^4R^2}
 \asymp \frac{N}{q_0^2x^{50\eta}H^2}.
\]
Thus $52=50+2$, with the additional $2\eta$ arising from the square of the
$x^\eta$ in the definition of $H$.  We retain this window throughout.

\subsection{The terminal minimum}

The more substantive issue occurs on page~18 and concerns
\begin{equation}\label{eq:delta-star-diagnosis}
 \Delta^*=\min\left\{
 \frac{N}{|\Lambda|x^{5\eta}},\Delta_1\right\}.
\end{equation}
The extension sketch gives lower bounds of the form
\[
 \frac{N}{|\Lambda|x^{5\eta}}\gg C,
 \qquad \Delta_1\gg C,
\]
and then sets $\Delta^*=\Delta_1$.  This conclusion does not follow: two
quantities being bounded below by the same third quantity gives no ordering
between them.  In particular, it does not establish
$\Delta_1\le N/(|\Lambda|x^{5\eta})$.

This is not merely a typographical matter, because reciprocal powers of
$\Delta^*$ occur in the final estimates.  Our repair is to preserve the
minimum and use the exact identity
\begin{equation}\label{eq:delta-branch-diagnosis}
 \frac1{\Delta_1\Delta^*}
 =\max\left\{
 \frac{|\Lambda|x^{5\eta}}{N\Delta_1},
 \frac1{\Delta_1^2}
 \right\}.
\end{equation}
Both branches are carried through the exponent ledger in
\cref{sec:closure}.

\subsection{Scholarly claim boundary}

The observations above do not invalidate the published smooth-modulus theorem
\cite{Stadlmann2025}; they concern only the larger prescribed-factor
extension in \cite{Stadlmann2026}.  Nor do they imply that the desired
extension is false.  Sections \ref{sec:terminal}--\ref{sec:outer} give a
complete replacement argument: the supported local descent is proved first,
and only then consumed by the outer Cartesian and dispersion closure.  The
attribution used below is therefore: the smooth-modulus q-van der Corput
architecture is due to Stadlmann, while the factor-interface transport to
$\DS$, the genuine-minimum branch analysis, the supported descent, and the
corrected outer Cartesian closure are contributions of the present paper.
